\documentclass[11pt]{article}

\usepackage[margin=25mm]{geometry}

\usepackage{amsmath,amssymb,fontspec,unicode-math}

\IfFontExistsTF{DejaVu Serif}{\setmainfont{DejaVu Serif}}{\setmainfont{Latin Modern Roman}}

\IfFontExistsTF{Latin Modern Math}{\setmathfont{Latin Modern Math}}{\setmathfont{DejaVu Math TeX Gyre}}

\usepackage{xurl,hyperref}

\setlength{\parindent}{0pt}\setlength{\parskip}{6pt}

\begin{document}

\title{DisCO: Reinforcing Large Reasoning Models with Discriminative Constrained Optimization}

\author{ELVIA preliminary AI analysis}\date{21 October 2026}\maketitle

Preliminary assessment: Major.\par

Independent scientific review remains pending. Author adoption has not been requested.

The latest paper and finite-group implementation retain difficulty-dependent inclusion and the false automatic-monotonicity inference from a lower bound. The historical L-ratio dispatcher failure is fixed in the current author release. Other bounded mathematical, manifest and reporting concerns persist; published training outcomes have not been reconstructed.

\section*{Review scope}

Claim-focused reading of the identified paper and author-released materials. The arguments below distinguish mathematical contradictions, version-specific implementation defects and unresolved empirical consequences. No author code was executed in this revision.

\section*{Problem 1: finite squared-hinge penalties and exact feasibility}

Source: Paper pp. 1-2, 6, Eqs. (8)-(9); its cited penalty-theory source arXiv:2410.03955v4, pp. 8-10, Assumptions 1-4 and Theorem 1.

Assessment: supported; impact: minor

\subsection*{Source claim (paraphrase)}

The discussion following Eq. (\(9\)) relates a finite squared-hinge KL penalty to the constrained KKT system, while the abstract/introduction describe trust-region enforcement.

\subsection*{Analysis}

At a feasible point, the gradient of \(\beta[D(\theta)-\delta]_+^2\) is zero. An active constrained optimum with a positive KKT multiplier can nevertheless have a nonzero objective gradient, so it need not be stationary for this finite penalty. For \(\pi_{\rm old}=(1/2,1/2)\), \(\pi_\theta=(\operatorname{sigmoid}\theta,1-\operatorname{sigmoid}\theta)\), \(J(\theta)=\theta\), and \(D(\theta)=\log\cosh(\theta/2)\), the constrained maximum is \(\theta_*=2\operatorname{arcosh}(e^\delta)\). At \(\delta=10^{-4}\) and \(\beta=1000\), \(\theta_*\simeq0.02828474\), whereas the penalized stationary maximum has \(\theta\simeq0.25371689\) and \(D\simeq0.00802504>\delta\). The bounded objective \(J(\theta)=2\tanh(\theta/2)\) similarly permits \(D\simeq0.00794101\).

\subsection*{Finding and scope}

Finite penalization alone does not imply exact hard feasibility. The cited penalty theorem concerns approximate KKT conditions with tolerance-dependent assumptions and schedules; that theorem is not refuted. Special instances can have exactly feasible penalized optima, so a literal existential claim that this can happen is also not refuted.

\subsection*{Proposed correction}

State approximate stationarity/feasibility and the required assumptions. If exact feasibility is intended, provide an appropriate constrained-update argument. For the scalar example, \(\theta_*=2\operatorname{arcosh}(e^\delta)\) and \(\lambda_*=2/\tanh(\theta_*/2)\).

Proposed correction; independent verification and author adoption remain unestablished.

\section*{Problem 2: finite groups retain difficulty-dependent selection}

Source: paper pages: 1 7 8 19; paper equations: 7 10; code: verl/verl/trainer/ppo/core\_algos.py:289-314,377-400,470-503,571-602 verl/verl/trainer/config/ppo\_trainer.yaml:159-166 scripts/train/run\_disco\_logL\_1.5b\_8k.sh

Assessment: supported; impact: major

\subsection*{Source claim (paraphrase)}

The abstract and Table 1 describe complete removal of difficulty bias. The released loss excludes groups containing only one reward class and divides by the original question count.

\subsection*{Analysis}

Let a question have success probability \(p\), and let \(K\sim\operatorname{Binomial}(n,p)\) be its number of correct responses. A mixed group is included with probability \(h_n(p)=1-p^n-(1-p)^n\). Since excluded groups contribute zero but the denominator remains the original question count, the expected finite-group contribution retains this difficulty-dependent factor. A question with one correct and one incorrect possible outcome already yields this construction for the full DRO objective. For \(n=8\), \(h_8(0.01)\simeq0.0772553\) and \(h_8(1/2)=0.9921875\); the relative inclusion weight is about \(0.0778636\). The cited configuration disables dynamic filtering, and the rejection-sampling flag defaults to false.

\subsection*{Finding and scope}

The finite fixed-group implementation does not completely remove difficulty-dependent question selection. The population decomposition correctly removes its explicit \(\sqrt{p(1-p)}\) factor, and the method can still provide useful practical reweighting.

\subsection*{Proposed correction}

Restrict the claim to the population normalization and available mixed groups, or specify a justified question-level sampling/inclusion correction. Dividing by the number of valid groups alone does not equalize original-question inclusion.

Proposed correction; independent verification and author adoption remain unestablished.

\section*{Problem 3: a lower bound does not ensure monotonic improvement}

Source: Current arXiv v5, Eqs.(7),(10), PDF pp.5,7, and paragraph following Eq.(10)

Assessment: supported; impact: major

\subsection*{Source claim (paraphrase)}

The paragraph following Eq. (\(10\)) infers that increasing the lower bound \(J_2\le J_1\) automatically increases \(J_1\).

\subsection*{Analysis}

The witness has one question and three one-token outputs with rewards \(r(o_1)=1\), \(r(o_2)=r(o_3)=0\). Under the fixed old policy \((0.4,0.3,0.3)\), the conditional positive law is a point mass on \(o_1\), and the conditional negative law assigns \(1/2\) to each of \(o_2,o_3\). With the paper’s log-likelihood score \(s_\theta(o_i)=\log\pi_i\), Eqs.(7) and (10) reduce exactly to \[J_1(\pi)=\log\pi_1-\tfrac12(\log\pi_2+\log\pi_3),\] \[J_2(\pi)=\log\pi_1-\tau\log\{(\pi_2^{1/\tau}+\pi_3^{1/\tau})/2\}.\] The trust-region quantity is \(D_{\rm KL}(\pi_{\rm old}\|\pi)=\sum_i\pi_{{\rm old},i}\log(\pi_{{\rm old},i}/\pi_i)\).

Write \(J_1=J_2+G\), where \(G\ge0\). An increase in \(J_2\) can be outweighed by a decrease in \(G\). A policy-level example uses reference probabilities \((0.4,0.3,0.3)\), temperature \(\tau=10\), and candidate policies \(A=(0.4,0.305,0.295)\) and \[B=(0.4000318,0.2999841,0.2999841)\]. Substitution in the paper's objectives gives \(J_2(B)-J_2(A)\simeq7.481517\times10^{-6}>0\), but \(J_1(B)-J_1(A)\simeq-6.409938\times10^{-6}<0\). Their KL values are approximately \(8.334491\times10^{-5}\) and \(2.106713\times10^{-9}\), both below \(\delta=10^{-4}\).

\subsection*{Finding and scope}

The asserted monotonic implication is false even within the stated KL region. The lower bound and the entropic DRO variational identity can remain correct; practical superiority is a separate empirical question.

\subsection*{Proposed correction}

Remove automatic monotonicity, or supply a touching-minorizer, gap-control or acceptance condition that actually controls \(J_1\).

Proposed correction; independent verification and author adoption remain unestablished.

\section*{Problem 4: Relative objective/penalty signs in Algorithm 1}

Source: paper pages: 19; algorithm lines: 9 10 11; code: verl/verl/trainer/ppo/core\_algos.py

Assessment: supported; impact: minor

\subsection*{Source claim (paraphrase)}

Algorithm 1 defines \(G_1\) as the objective gradient and \(G_2\) as the positive penalty gradient, then supplies \(G_1+G_2\) to the optimizer.

\subsection*{Analysis}

Descent on \(-J_2+P\) requires \(-G_1+G_2\); ascent requires \(G_1-G_2\). Adding both gradients gives the wrong relative sign under either convention. The author-released loss minimizes the negative objective plus the positive penalty, so its relative sign is correct.

\subsection*{Finding and scope}

Following pseudocode can give wrong direction; this does not establish wrong public loss or wrong reported training. Actual public objective/penalty relative sign is correct.

\subsection*{Proposed correction}

For a descent optimizer, write \(-G_1+G_2\) and state the convention.

Proposed correction; independent verification and author adoption remain unestablished.

\section*{Problem 5: The historical L-ratio dispatcher defect is fixed}

Source: Author repository Optimization-AI/DisCO, commit 8bdd94eb49a58cd09e68af5392302ba83acf63b2; verl/verl/workers/actor/dp\_actor.py, lines 286–296; verl/verl/trainer/ppo/core\_algos.py, compute\_policy\_loss\_disco\_Lratio

Assessment: not\_issue; impact: no\_issue\_found

\subsection*{Source claim (paraphrase)}

The original supplemental release routed the full L-ratio branch to a basic loss function that did not accept its tau argument.

\subsection*{Analysis}

In current author commit 8bdd94eb49a58cd09e68af5392302ba83acf63b2, dp\_actor.py, lines 286–296, calls compute\_policy\_loss\_disco\_Lratio, whose signature accepts tau. Reading the current call and signature clears the particular dispatch failure. The older supplemental call remains a historical version fact, and no author program was executed in this check.

\subsection*{Finding and scope}

This specific failure is absent from the current author release. Historical experiment provenance and adoption of any external audit remain unestablished.

\section*{Problem 6: Printed and implemented penalty gradients differ by a factor of two}

Source: paper pages: 6 19; paper equations: 9; code: verl/verl/trainer/ppo/core\_algos.py:318-321,406-409,506-509,604-607

Assessment: supported; impact: minor

\subsection*{Source claim (paraphrase)}

Eq. (\(9\)) uses \(\beta[D-\delta]_+^2\), whose gradient is \(2\beta[D-\delta]_+\nabla D\).

\subsection*{Analysis}

In each cited implementation branch, the detached coefficient gives only \(\beta[D-\delta]_+\nabla D\). Detachment removes the coefficient derivative, making this the gradient of \((\beta/2)[D-\delta]_+^2\).

\subsection*{Finding and scope}

Numerical beta setting is not an exact Eq \(9\) reproduction; retuning can absorb factor, so less severe than sign error. Penalty direction is correct; beta is tunable.

\subsection*{Proposed correction}

Use a consistent \(\beta/2\) convention or double the detached coefficient. Differentiating \(\beta\max(D-\delta,0)^2\) directly gives the printed gradient.

Proposed correction; independent verification and author adoption remain unestablished.

\section*{Problem 7: AMC benchmark year differs from the released evaluation manifest}

Source: paper pages: 7 8 9 18; paper tables: 2 3 4 5; data: deepscaler/deepscaler/data/test/amc.json datasets/deepscaler/data/amc.parquet; code: deepscaler/deepscaler/data/utils.py scripts/data/deepscaler\_dataset.py scripts/eval/eval\_model\_8k.sh

Assessment: supported; impact: minor

\subsection*{Source claim (paraphrase)}

Task setting/table headers identify AMC 2023.

\subsection*{Analysis}

AMC JSON has \(83\) rows: \(43\) source URLs dated \(2022\) and \(40\) dated 2023. Parquet has same \(83\) problems in same order. Loader/preprocessor/eval path contains no year filter; relevant files unchanged in pinned GitHub.

\subsection*{Finding and scope}

Supplied pipeline differs from named benchmark; external comparisons are ambiguous. Artifacts alone do not establish whether tables used \(83\) questions or undocumented \(40\)-question subset.

\subsection*{Proposed correction}

Publish actual table evaluation manifest. Supply/filter \(40\)-question \(2023\) split or relabel mixed-year results.

Proposed correction; independent verification and author adoption remain unestablished.

\section*{Problem 8: MATH500 overlaps and duplicate training pairs}

Source: paper pages: 7; data: deepscaler/deepscaler/data/train/deepscaler.json deepscaler/deepscaler/data/test/math.json datasets/deepscaler/data/train.parquet datasets/deepscaler/data/math.parquet

Assessment: supported; impact: minor

\subsection*{Source claim (paraphrase)}

Paper describes about \(40.3\) thousand unique pairs; delivered train/test files should identify benchmark separation.

\subsection*{Analysis}

The author-released MATH test indices \(80,285,412\) match training indices \(7202,7141,6976\) after Unicode NFKC normalization, whitespace normalization and case folding; the corresponding answers \(501,440,406\) agree. The first pair is also identical before normalization. JSON and Parquet problem orders agree. There are \(40\,315\) training rows but only \(39\,616\) raw unique problem-answer pairs, \(39\,387\) raw unique problems and \(39\,097\) normalized unique problems. The other five released test sets have no exact normalized match under this narrow comparison.

\subsection*{Finding and scope}

These are concrete overlaps in the delivered data, with no semantic/fuzzy matching claim. Even attributing all three MATH answers to overlap changes an individual six-set macro-score by at most \(0.1\) percentage point; this does not explain the full reported improvement.

\subsection*{Proposed correction}

Correct counts; publish deduplicated/decontaminated manifests. Report sensitivity excluding \(3\) affected MATH questions; use clean future training release. Retraining is required to estimate causal effect of changing training data.

Proposed correction; independent verification and author adoption remain unestablished.

\section*{Problem 9: Checkpoint selection and uncertainty: unresolved empirical scope}

Source: paper pages: 7 8 9 10 24 25; code: scripts/train/run\_disco\_logL\_1.5b\_8k.sh scripts/eval/eval\_model\_8k.sh

Assessment: supported; impact: minor

\subsection*{Source claim (paraphrase)}

Paper reports best performance and significant superiority without a complete selection manifest or uncertainty intervals.

\subsection*{Analysis}

The paper evaluates every \(200\) steps and reports best performance, without a fully specified independent checkpoint-selection rule. Its checklist openly reports no error bars. Averaging \(16\) correctness samples estimates pass@\(1\), rather than pass@\(16\), and does not measure training-seed variability. Dynamic sampling is disabled for all compared methods, which is disclosed. The fixed \(0.001\) KL ablation does not establish failure of every tuned regularizer.

\subsection*{Finding and scope}

Multiple models, six benchmarks and a second training dataset provide useful reported evidence. Checkpoint attribution and statistical robustness remain unresolved; no corrected effect size is supplied.

\subsection*{Proposed correction}

Publish table/checkpoint/seed/config/dataset manifests and selection protocol. Compute paired question-level intervals from saved outputs; estimate training-seed uncertainty separately. Scope DAPO and stability comparisons to tested conditions.

Proposed correction; independent verification and author adoption remain unestablished.

\section*{Problem 10: Proposition 1's sign-dependent monotonicity}

Source: paper pages: 4 18 19; paper equations: 3 4 5

Assessment: supported; impact: minor

\subsection*{Source claim (paraphrase)}

Proposition 1 assumes f nondecreasing in x and remark says assumptions hold for GRPO.

\subsection*{Analysis}

For a negative reward \(y=-1\), the clipped branch is \(f(x,-1)=-\max(x,1-\varepsilon)\), which is nonincreasing. At \(\varepsilon=0.2\), \(f(1,-1)=-1\) and \(f(1.1,-1)=-1.1\), contradicting the claimed nondecreasing condition. The population decomposition itself does not require that extra condition.

\subsection*{Finding and scope}

Correct the sign-dependent assumption; the decomposition identity is not refuted.

\subsection*{Proposed correction}

Remove extra monotonicity of f for all y; retain branchwise assumptions or state sign-dependent monotonicity.

Proposed correction; independent verification and author adoption remain unestablished.

\section*{Problem 11: Dr. GRPO's relative question weighting}

Source: paper pages: 4 5 20 21; paper figures: 1; paper equations: 19

Assessment: supported; impact: minor

\subsection*{Source claim (paraphrase)}

The text interprets \(p(1-p)\) as mitigating question imbalance relative to \(\sqrt{p(1-p)}\).

\subsection*{Analysis}

After normalizing both weights to one at \(p=1/2\), \(r_{\rm GRPO}(p)=2\sqrt{p(1-p)}\) and \(r_{\rm Dr}(p)=4p(1-p)=r_{\rm GRPO}(p)^2\). At \(p=0.01\), these are approximately \(0.1989975\) and \(0.0396\). Under the displayed relative-weight criterion, the latter attenuates extreme questions more strongly.

\subsection*{Finding and scope}

The interpretation of the relative weights is reversed. This does not determine empirical performance, and an overall gradient-scale change can be absorbed into the learning rate.

\subsection*{Proposed correction}

Distinguish absolute scale from relative weighting and other definitions of bias; correct mitigation sentence under displayed criterion.

Proposed correction; independent verification and author adoption remain unestablished.

\section*{Limits}

Independent scientific adjudication remains pending.\par

No author code, model training or full benchmark evaluation was executed in this revision.\par

Version-specific source defects do not establish which code produced the published experiments.\par

Corrected benchmark scores, empirical effect sizes and author adoption remain unverified unless explicitly stated.\par

\section*{Source identity}

arXiv:2505.12366v5, revised 6 January 2026

\url{https://arxiv.org/pdf/2505.12366v5}

Reviewed PDF SHA-256: \nolinkurl{947fc7d781d73ca93dd10e06e6aba99621f845d6a503a896379a7fb6a7b53bbb}

This latest arXiv revision is the same DisCO work as the original conference paper, confirmed by its title, five authors and operative claims. The exact current primary PDF inspected is identified by its SHA-256. Earlier immutable source identity remains in the predecessor record. The download timestamp is recorded from this fresh primary-source retrieval.

\end{document}
